% Figure: the p-test for improper integrals. Left: tails 1/x^p on
% [1, infinity), converging for p > 1. Right: near-origin singularities
% 1/x^p on (0, 1], converging for p < 1. The boundary p = 1 is the
% divergent case in both.
\begin{figure}[htbp]
\centering
\begin{tikzpicture}

\begin{axis}[
  name=tail,
  width=7cm, height=5.2cm,
  xlabel={$x$}, ylabel={$x^{-p}$},
  xmin=1, xmax=6, ymin=0, ymax=1.05,
  domain=1:6, samples=150,
  axis lines=left, grid=both,
  major grid style={engblack!12}, font=\footnotesize,
  title={Tail at $\infty$: converges for $p>1$},
  legend pos=north east, legend style={font=\scriptsize},
]
\addplot[engred, very thick] {x^(-0.5)}; \addlegendentry{$p=0.5$ (div.)}
\addplot[engblack, very thick] {x^(-1)}; \addlegendentry{$p=1$ (div.)}
\addplot[engblue, very thick] {x^(-2)}; \addlegendentry{$p=2$ (conv.)}
\end{axis}

\begin{axis}[
  at={(tail.east)}, anchor=west, xshift=12pt,
  width=7cm, height=5.2cm,
  xlabel={$x$}, ylabel={},
  xmin=0, xmax=1, ymin=0, ymax=6,
  domain=0.03:1, samples=200,
  axis lines=left, grid=both,
  major grid style={engblack!12}, font=\footnotesize,
  title={Origin: converges for $p<1$},
  legend pos=north east, legend style={font=\scriptsize},
]
\addplot[engblue, very thick] {x^(-0.5)}; \addlegendentry{$p=0.5$ (conv.)}
\addplot[engblack, very thick] {x^(-1)}; \addlegendentry{$p=1$ (div.)}
\addplot[engred, very thick] {x^(-1.5)}; \addlegendentry{$p=1.5$ (div.)}
\end{axis}

\end{tikzpicture}
\caption[The p-test for improper integrals]{The two power-law tests
for improper integrals, mirror images of each other. Left: on the
infinite tail $[1, \infty)$ the integral of $x^{-p}$ converges when
the integrand decays fast enough, that is $p > 1$, and diverges for
$p \leq 1$. Right: near a singularity at the origin on $(0, 1]$ the
integral of $x^{-p}$ converges when the integrand blows up slowly
enough, that is $p < 1$, and diverges for $p \geq 1$. The boundary
$p = 1$ diverges in both settings, the logarithmic borderline that an
engineer commits to memory as the dividing line for tail bounds on
densities and for endpoint singularities in field integrals.}
\label{fig:vol02:ch06:improper-ptest}
\end{figure}
